\section*{3. 与线束相伴的爆破}

\subsection*{3.1.}
设 $X$ 为代数闭域 $k$ 上的 $n$ 维不可约光滑射影概形, $D=\{H_\tau\}_{\tau\in\mathbb P^1}$ 为一个超曲面线束, 其轴与 $X$ 横截相交于 $\Delta$. 我们考虑 $X$ 沿 $\Delta$ 的爆破, 现在将其显式描述如下.\footnote{印本在本句中用 $u$ 表示维数. 第 3.2 节随即用 $n$ 表示同一维数, 而 (3.2.3)--(3.2.4) 中的符号也由它决定; 此处采用后一记号.}

用方程表示, 设 $(\lambda,\mu)$ 为 $\mathbb P^1$ 上的齐次坐标, $F$ 和 $G$ 为两个线性型, 使得 $F=G=0$ 定义该线束的轴. 该线束写为
\begin{equation}\tag{3.1.1}
\lambda F-\mu G.
\end{equation}
按定义, 子概形 $\Delta$ 是 $X$ 与线性子簇 $G=F=0$ 的交. 由 2.1.6,
\begin{equation}\tag{3.1.2}
\widetilde X
=\{(x,(\lambda,\mu))\in X\times\mathbb P^1
       \mid \lambda F-\mu G=0\};
\end{equation}
它是映射
\begin{equation}\tag{3.1.3}
\begin{tikzcd}[row sep=large,column sep=huge]
X-\Delta \arrow[r] & \mathbb P^1 \\
x \arrow[r] & (G(x),F(x)).
\end{tikzcd}
\end{equation}
的图像在 $X\times\mathbb P^1$ 中的概形论闭包. $\widetilde X$ 的重要优点在于它在 $k$ 上为射影概形 (因为嵌入 $X\times\mathbb P^1$ 中) 且光滑 (因为 $\Delta$ 光滑), 并有到 $\mathbb P^1$ 的态射 $\rho$, 即
\begin{equation}\tag{3.1.4}
\rho:\widetilde X\longrightarrow X\times\mathbb P^1
\xrightarrow{\operatorname{pr}_2}\mathbb P^1,
\end{equation}
其纤维为该线束中各成员在 $X$ 上截得的超平面截面:
\[
\rho^{-1}(t)=X_t\overset{\mathrm{def}}{=}X\cdot H_t.
\]

\subsection*{3.2. 对偶簇的次数: 拓扑方法}

再假设 3.1 中的线束满足
\begin{equation}\tag{3.2.1}
X_t\text{ 光滑, 当 }t\text{ 遍历下列空间的某个非空开子集时: }
\mathbb P^1,
\end{equation}
以及
\begin{equation}\tag{3.2.2}
\text{若 $X_t$ 奇异, 则它的奇点仅为非退化二次奇点 (XVII 1.1).}
\end{equation}
特别地, 若 $\operatorname{char}k\ne2$ 或 $X$ 的维数为偶数, 则\underline{Lefschetz 线束} (XVII 2.2) 满足这些条件.

设 $\eta$ 为 $\mathbb P^1$ 的泛点, $\bar\eta$ 为其上的几何点, $n=\dim X$, $\chi$ 为以 $\mathbb Q_\ell$ 为系数的 Euler--Poincar\'e 示性数. 将 Euler--Poincar\'e 公式 XVI 2.1 应用于 $\rho$, 得
\begin{equation}\tag{3.2.3}
\chi(\widetilde X)
=\chi(\mathbb P^1)\,\chi(X_{\bar\eta})
 +(-1)^n\sum_{\substack{\text{各纤维的}\\
                         \text{奇点}}}1,
\end{equation}
由 (2.2.2), 可将它改写为
\begin{equation}\tag{3.2.4}
(-1)^n
\left[\begin{array}{c}
\text{奇点}\\[3pt]
\text{总数}
\end{array}\right]
=\chi(X)+\chi(\Delta)-2\chi(X_{\bar\eta}).
\end{equation}
注意, 量 $\chi(X)+\chi(\Delta)-2\chi(X_{\bar\eta})$ 与线束的具体选取无关. 它只依赖于 $X$ 在某个 $\mathbb P^r$ 中的嵌入, 实际上只依赖于 $\underline{\mathcal O}_X(1)$ 的 Chern 类; 见 XVII 5.7.5 和 5.7.6.

\begin{editorialnote}
底本此处将 Lefschetz 线束的定义指向 ``XVII 2.21, 但第 XVII 讲的定义在 2.2, 不存在 2.21. 此处据定义内容订正为 XVII 2.2.
\end{editorialnote}

\subsubsection*{3.2.5.}
我们在一个具体情形中将其写出. 取 $X=\mathbb P^2$, 设 $\{H_\tau\}$ 为满足 3.1、(3.2.1) 和 (3.2.2) 的 $d$ 次平面曲线线束. 则
\begin{equation}\tag{3.2.6}
\begin{array}{ll}
\chi(X)=3 & \text{(由 1.2)},\\[4pt]
\chi(X_{\bar\eta})=-d(d-3) & \text{(见 SGA 5 VII 7.3)},\\[4pt]
\chi(\Delta)=d^2,&
\end{array}
\end{equation}
从而
\begin{equation}\tag{3.2.7}
\text{奇点总数为 }3(d-1)^2.
\end{equation}
取特殊情形 $d=3$, 并考虑椭圆曲线族 (特征不等于 $3$)
\begin{equation}\tag{3.2.8}
X^3+Y^3+Z^3+3\mu XYZ.
\end{equation}
四个奇异纤维分别位于 $\mu^3=-1$ 和 $\mu=\infty$ 之上,\footnote{印本称有限参数为三次单位根. 对 (3.2.8) 中印出的加号, Hesse 分解给出 $\mu^3=-1$; $+1$ 的三次根对应于系数为 $-3\mu$ 的多项式.}
且每个奇异纤维都是三条直线之并, 如下所示:
\begin{equation}\tag{3.2.9}
\begin{tikzcd}[row sep=4.5em,column sep=4.5em]
{} & {} \arrow[dl,dash,shorten <=-9mm,shorten >=-9mm]
        \arrow[dr,dash,shorten <=-9mm,shorten >=-9mm] & {} & {} \\
{} \arrow[rr,dash,shorten <=-11mm,shorten >=-11mm] & {} & {} & {,}
\end{tikzcd}
\end{equation}
它恰有三个奇点, 全都是非退化二次奇点. 总计得到十二个奇点, 在本例中验证了 (3.2.7).

\subsubsection*{命题 3.2.10.}
\underline{设 $X\hookrightarrow\mathbb P^r$ 为 Lefschetz 嵌入}
(XVII 2.3), \underline{并假设 $n=\dim X$ 为偶数, 或 $\operatorname{char}k\ne2$. 则每个 Lefschetz 线束看作 $\check{\mathbb P}^{r}$ 中的直线时, 与 $\check X$ 相交于}
\begin{equation}\tag{3.2.11}
(-1)^n\bigl[\chi(X)+\chi(\Delta)-2\chi(X_{\bar\eta})\bigr]
\end{equation}
\underline{个点, 且与之横截相交} (EGA IV 17.13.7).

\underline{证明.}
由假设, Lefschetz 线束的每个纤维的奇点都是非退化二次奇点. 因而公式 (3.2.4) 表明, 线束 $D$ 中奇异纤维的个数 (每个这样的纤维恰含一个奇点) 就是 (3.2.11) 中的数. 而纤维 $X_\tau$ 奇异当且仅当 $\tau\in\check X\cap D$.

为说明相交横截, 注意 (3.2.11) \underline{不依赖于}具体的 Lefschetz 线束. 由于一般的线束是 Lefschetz 线束, (3.2.11) 就是 $\check X$ 的次数; 即它作为 $\check{\mathbb P}^{r}$ 中超曲面的次数, 约定在 $\dim\check X\le r-2$ 时此次数为 $0$. 若直线 $D$ 与一个 $d>0$ 次超曲面相交于 $d$ 个不同的点 $x_i$ ($1\le i\le d$), 则此交必为横截交. 事实上, 局部相交重数
\[
\mu_i=\operatorname{length}_{x_i}(D\cap\check X)
\]
至少为 $1$, 而 $\sum_{i=1}^d\mu_i=d$. 因而对每个 $i$ 都有 $\mu_i=1$; 等价地, $D\cap\check X$ 光滑, 这正是横截性.\footnote{印本以 $\operatorname{rg}$ 表示局部长度; 标准的局部长度记号以及所述相交重数论证都要求使用 $\operatorname{lg}$ (或等价的 $\operatorname{length}$).}

\subsubsection*{注 3.2.12.}
对于已经假设与 $\check X$ 横截相交的 Lefschetz 线束 $D$ (充分一般的 $D$ 总满足此条件), 结果 3.2.10 也可看作 XVII 5.7.2 与 XVII 3.5 中 $\deg\varphi=1$ 这一事实的结合. 当然, 此处借助 Euler--Poincar\'e 公式 XVI 2.1 证明的\underline{每个} Lefschetz 线束 $D$ 都与 $\check X$ 横截相交这一事实, 也可利用 $D$ 的轴与 $X$ 横截相交这一事实, 直接而初等地验证.
